\documentclass[12pt]{article}
\usepackage[utf8]{inputenc}
\usepackage[T1]{fontenc}
\usepackage{hyperref}
\title{Warembori Language Resources for K-12 Curriculum Development}
\author{Compiled by Claude Sonnet 5.5 (claude-sonnet-5-5) for Andr\'as Kornai}
\date{}
\begin{document}
\maketitle

\section{Introduction}
Warembori (ISO 639-3 wsa, Glottocode ware1253) is a moribund language of the north coast of Papua, Indonesia. It is spoken in Warembori village, Mamberamo Hilir District, Mamberamo Raya Regency. The sources I found give about 600 speakers as of 1998. Its affiliation is debated. Donohue groups it with Yoke as the Lower Mamberamo family, while Ross and later researchers have treated it as a Papuanised Austronesian language. It is not known to be taught in schools.

\section{Primary Educational Resources}
No primers, textbooks, school programmes or orthography materials for Warembori were found.

\section{Linguistic and Reference Materials}
\begin{itemize}
\item Wikipedia, Warembori language: overview and classification. \url{https://en.wikipedia.org/wiki/Warembori_language}
\item Wikipedia, Lower Mamberamo languages: family context. \url{https://en.wikipedia.org/wiki/Lower_Mamberamo_languages}
\item Mark Donohue, a short note on Warembori hosted at the Max Planck Institute EVA. It discusses verb agreement similar to Kwerba, Austronesian-like plural marking, and the question of classification. \url{https://www.eva.mpg.de/lingua/tools-at-lingboard/pdf/donohue_waremborinote.pdf}
\item Endangered Languages catalogue entry for Warembori. \url{https://endangeredlanguages.com/es/node/88791}
\item Endangered Languages catalogue sample page, which lists a 1999 short introduction to the grammar of Warembori. I did not open it. \url{https://www.endangeredlanguages.com/lang/3322/samples/10754?hl=de}
\item Rosetta Panglossia entry for the code wsa, with alternate names. \url{https://rosettapanglossia.longnow.org/index.php?title=Wsa}
\end{itemize}

\section{Religious and Community Resources}
No scripture translations or community organisations for Warembori were found.

\section{Digital and Audio Resources}
No audio, video, corpus or software resources were found.

\section{Curriculum Development Context}
Documentation is thin: only short sketches and notes, with no full reference grammar or dictionary found. A curriculum would first need community-led documentation with the speakers, who are mainly adults. Indonesian would be the likely language of schooling. The resources above are suitable mainly as linguistic background for such work.
\end{document}